\section{总结与延伸}

\subsection{全课知识图谱}

本节课覆盖了从 CNN 组件到完整训练流程的全部内容：

\begin{center}
\begin{tikzpicture}[node distance=0.7cm, every node/.style={draw, rounded corners, minimum width=3.5cm, minimum height=0.6cm, font=\small}]
\node (layers) {CNN 层组件};
\node (norm) [below=of layers] {归一化与正则化};
\node (act) [below=of norm] {激活函数};
\node (arch) [below=of act] {经典架构（VGG, ResNet）};
\node (train) [below=of arch] {训练技巧};
\draw[->, thick] (layers) -- (norm);
\draw[->, thick] (norm) -- (act);
\draw[->, thick] (act) -- (arch);
\draw[->, thick] (arch) -- (train);
\node[draw=none, right=0.8cm of layers, text width=5.5cm, font=\scriptsize] {卷积、池化、全连接};
\node[draw=none, right=0.8cm of norm, text width=5.5cm, font=\scriptsize] {Layer Norm、Batch Norm、Dropout};
\node[draw=none, right=0.8cm of act, text width=5.5cm, font=\scriptsize] {Sigmoid $\to$ ReLU $\to$ GELU};
\node[draw=none, right=0.8cm of arch, text width=5.5cm, font=\scriptsize] {小滤波器、残差连接、深度可扩展};
\node[draw=none, right=0.8cm of train, text width=5.5cm, font=\scriptsize] {初始化、学习率调度、数据增强、迁移学习};
\end{tikzpicture}
\end{center}

\subsection{关键 Takeaways}

\begin{importantbox}{七条核心原则}
\begin{enumerate}
    \item \textbf{归一化稳定训练}：Layer Norm（Transformer）或 Batch Norm（CNN）让激活值分布可控
    \item \textbf{Dropout 防止过拟合}：通过随机丢弃强制学习冗余表示
    \item \textbf{GELU 取代 ReLU}：平滑且无死神经元，已成为 Transformer 标配
    \item \textbf{小滤波器更好}：$3 \times 3$ 卷积参数少、非线性强，VGG 证明了这一点
    \item \textbf{残差连接使深度成为可能}：学习增量 $F(x)$ 比学习完整映射 $H(x)$ 容易得多
    \item \textbf{正确初始化至关重要}：Kaiming 初始化根据维度和激活函数校准权重标准差
    \item \textbf{训练是系统工程}：学习率调度、数据增强、迁移学习缺一不可
\end{enumerate}
\end{importantbox}

\subsection{拓展阅读}

\begin{itemize}
    \item Simonyan \& Zisserman, ``Very Deep Convolutional Networks'' (VGGNet, 2014)
    \item He et al., ``Deep Residual Learning for Image Recognition'' (ResNet, 2015)
    \item He et al., ``Delving Deep into Rectifiers'' (Kaiming 初始化, 2015)
    \item Ioffe \& Szegedy, ``Batch Normalization'' (2015)
    \item Ba et al., ``Layer Normalization'' (2016)
    \item Hendrycks \& Gimpel, ``Gaussian Error Linear Units (GELUs)'' (2016)
    \item Srivastava et al., ``Dropout: A Simple Way to Prevent Neural Networks from Overfitting'' (2014)
\end{itemize}

\end{document}
